\section{Nonsymmetric cones: barrier parameters and evaluation costs}
\label{sec:nonsymmetric-barriers}
The dimension $m$ of a cone and its contact capacity (at most $m-2$)
do not determine its best barrier parameter. We bound that parameter for the norm and perspective cones of
Section~\ref{sec:lp-power}. Throughout, $\nu_{\rm LH}(C)$ denotes the
infimum of the parameters of logarithmically homogeneous
self-concordant (LHSC) barriers on $C$; on a product, these barriers
may couple the factors. The cross-ratio arguments below require
logarithmic homogeneity.

\subsection{An additive lower bound for fixed cone dictionaries}
For $1<p<\infty$, put $c=\min\{1/p,1-1/p\}$. If $p\ne2$, let
$\gamma\in(0,1)$ solve
\begin{equation}\label{eq:hildebrand-root}
 (1-2c)\gamma^{1-c}+(1-c)\gamma^{1-2c}=c,
 \qquad h(p)=1+\frac{1+\gamma^c}{1+\gamma^{1-c}}.
\end{equation}
Set $h(2)=2$. The left side of the root equation is strictly increasing
from zero to $2-3c>c$, so the root exists and is unique. Rearranging gives
\begin{equation}\label{eq:hildebrand-excess}
 h(p)-2=\frac{1-2c}{1-c}\gamma^c.
\end{equation}
Consequently $2<h(p)<3$ for $p\ne2$, and $h(p)=h(q)$ when
$1/p+1/q=1$. The lower bound $h(p)$ for a single three-dimensional
power cone or norm cone is due to
\citet[Corollaries~7.1--7.2]{Hildebrand2013}.

For a proper cone $C$, choose a compact affine base, a basis
$e_1,\ldots,e_m$ of boundary points of this base, and a point $x$ in
the relative interior of their simplex. On the projective line through
$e_a$ and $x$, let $y_a$ be the other boundary point of the base and
$z_a$ the intersection with the span of the other basis vectors. Write
$q_a=(e_a,x;y_a,z_a)$, using the convention
$(a,b;c,d)=(a-c)(b-d)/((a-d)(b-c))$ and continuity at infinity.
Hildebrand's Theorem~5.4 states, in particular, that
\begin{equation}\label{eq:positive-crossratio}
 \sum_aq_a\geq2\quad\Longrightarrow\quad
 \nu_{\rm LH}(C)\geq\sum_aq_a.
\end{equation}

\begin{lemma}[Cross-ratio lower bound for products]
\label{lem:positive-product}
Suppose each proper cone $C_i$ has a configuration as above with
cross-ratio sum $S_i\geq2$. Then
$\nu_{\rm LH}(\prod_i C_i)\geq\sum_i S_i$.
\end{lemma}
\begin{proof}
Normalize each local base by $\ell_i=1$ and the product base by
$\sum_i\ell_i=1$. Embed each $e_{ia}$ into its own direct summand, and
choose $X=(\alpha_i x_i)_i$, where $\alpha_i>0$ and $\sum_i\alpha_i=1$.
These embedded vectors form a boundary basis, and $X$ lies in the
interior of their simplex. Along the line from an embedded $e_{ia}$ through $X$
and beyond, every other block remains a positive multiple of $x_j$.
The point where the line leaves the base beyond $X$, and its
intersection with the span of the other basis vectors, are therefore
determined by block $i$.
Projection onto block $i$ induces a projective isomorphism on this line,
sending the four defining points to their local counterparts. Thus its
cross-ratio is $q_{ia}$, unchanged by $\alpha_i$. Applying
\eqref{eq:positive-crossratio} to the product gives the result.
\end{proof}

\begin{theorem}[Lower bounds for products of nonsymmetric cones]
\label{thm:nonsymmetric-premium}
For $m_i\geq3$ and $r_i\geq1$,
\begin{align}
 \nu_{\rm LH}\left(\prod_{i=1}^k K_{p,m_i}\right)&\geq k h(p),
 \label{eq:porder-premium}\\
 \nu_{\rm LH}\left(\prod_{i=1}^k\mathcal P_{p,r_i}\right)&\geq k h(p).
 \label{eq:perspective-premium}
\end{align}
At $p=2$ both values of $\nu_{\rm LH}$ equal $2k$.
\end{theorem}
\begin{proof}
For $1<p<2$, Hildebrand's norm-cone configurations have positive
cross-ratio sums tending to $h(p)>2$. Apply
Lemma~\ref{lem:positive-product} and take the limit. For $p>2$, the
signed Legendre transform
$F_*(s)=\sup_{x\in\intr C}\{-\ip sx-F(x)\}$ preserves the LHSC
parameter and has domain $\intr C^*$. Since
$K_{p,3}^*=K_{q,3}$, the preceding case applied to the dual product
proves the assertion. This dual step is necessary: for $p>2$,
Hildebrand's configurations have $\sum_aq_a<2$, so
Lemma~\ref{lem:positive-product} does not apply to them directly.
In a larger norm cone, set all but two norm coordinates to zero. This
linear section meets the ambient interior and is $K_{p,3}$, so barrier
restriction proves \eqref{eq:porder-premium}.

In $\mathcal P_{p,r_i}$ retain one output coordinate. The resulting
interior-meeting section is the scalar power cone
$u^{1/p}v^{1/q}\geq|w|$. For $p\geq2$, Hildebrand's power-cone
configurations have $\sum_aq_a\geq2$; for $p<2$, swap $u$ and $v$. Lemma~\ref{lem:positive-product} proves
\eqref{eq:perspective-premium}. At $p=2$ every displayed cone is a
Lorentz cone; the orthant lower bound and the sum of the standard
Lorentz barriers both give $2k$.
\end{proof}

The theorem gives additive lower bounds for a \emph{fixed} product,
without assuming a separable barrier. With
Theorem~\ref{thm:lp-granularity}, it shows that every lift of $B_p^N$
by norm-cone factors of dimensions between three and $d$ has an ambient
cone with $\nu_{\rm LH}\geq h(p)\lceil(N-1)/(d-2)\rceil$.
At the minimum dimension, Theorem~\ref{thm:minimal-dimension-rigidity}
gives a dictionary-independent statement: in any lift with $N+1$ cone
coordinates, the lifting cone is isomorphic to $K_{p,N+1}$ and so has
$\nu_{\rm LH}\geq h(p)$.

The bound $h(p)$ per factor does not extend to arbitrary factors
merely because the target is a non-Euclidean ball. For rational $p$, scalar power
inequalities have finite Lorentz lifts \citep{Wang2024}; applying these
to $|x_j|^p\leq u_j$, $\sum_j u_j=1$, gives an exact Lorentz-product
lift of $B_p^N$. Each of its factors has optimal parameter two.
This does not determine the best parameter over all dictionaries with
a dimension cap, because this lift may use more than the minimum
number of factors.

\subsection{Grouped perspective cones: duality and parameter bounds}
The perspective cones are known constructions
\citep{GlineurTerlaky2004,Chares2009}. Their exact dual is
\begin{equation}\label{eq:perspective-dual}
 \mathcal P_{p,r}^*=
 \{(a,b,z):a,b\geq0,\quad
 a^{1/p}b^{1/q}\geq p^{-1/p}q^{-1/q}\norm z_q\}.
\end{equation}
Indeed, writing $w=vs$ reduces dual membership to
$a\norm s_p^p+b+\ip zs\geq0$ for every $s$; the conjugate of
$a\norm s_p^p$ is $\norm z_q^q/[q(ap)^{q-1}]$ for $a>0$, with the
boundary cases obtained by closure. Thus the map
$(a,b,z)\mapsto(b,a,p^{-1/p}q^{-1/q}z)$ identifies the dual with
$\mathcal P_{q,r}$. The section $u=v=t$ also gives
\begin{equation}\label{eq:perspective-section-comparison}
 \nu_{\rm LH}(\mathcal P_{p,r})\geq\nu_{\rm LH}(K_{p,r+1}).
\end{equation}
The same comparison holds for products, by simultaneous restriction.
A dimension-$m$ proper cone admits an $m$-LHSC barrier, for example its
canonical barrier \citep{HildebrandCanonical2014}. Therefore
\begin{equation}\label{eq:perspective-interval}
 kh(p)\leq\nu_{\rm LH}\left(\prod_i\mathcal P_{p,r_i}\right)
 \leq\sum_i(r_i+2).
\end{equation}
We do not claim that either bound is exact for $p\ne2$.

We next give an explicit barrier with computable derivatives but a
larger parameter. Put
$\alpha=1/p$, $\beta=1/q$. Chares's scalar power barrier is
\begin{equation}\label{eq:chares-scalar-power}
 \phi_\alpha(a,v,z)
 =-\log(a^{2\alpha}v^{2\beta}-z^2)-\beta\log a-\alpha\log v.
\end{equation}
It has parameter three \citep[Theorem~3.1.1]{Chares2009}. Since
$\mathcal P_{p,r}$ is represented by
$\sum_i a_i=u$, $(a_i,v,w_i)\in\mathcal P_{p,1}$, define
\begin{equation}\label{eq:grouped-implicit-barrier}
 \Psi_{p,r}(u,v,w)=
 \min_{\sum_i a_i=u}\sum_i\phi_\alpha(a_i,v,w_i).
\end{equation}
For an interior point the fiber has a nonempty interior and lies in
$0<a_i<u$. The objective is strictly convex and diverges at its
boundary. The partial-minimization theorem for bounded fibers in
Section~\ref{sec:concrete-barriers} therefore gives a nondegenerate
self-concordant barrier. Scaling the unique minimizer proves logarithmic
homogeneity with parameter $3r$. This construction and the use of partial
minimization are due to Chares; the formulas below give its
evaluation cost.

At the minimizer, write $\phi_a(a_i,v,w_i)+\lambda=0$ and
$\sum_i a_i=u$. For a direction $(\dot u,\dot v,\dot w)$ set
\[
 A_i=\phi_{aa}>0,\qquad
 g_i=\phi_{av}\dot v+\phi_{aw}\dot w_i,\qquad S=\sum_i A_i^{-1}.
\]
Differentiation gives
\begin{equation}\label{eq:grouped-hessian-oracle}
 \dot\lambda=-\frac{\dot u+\sum_i g_i/A_i}{S},\qquad
 \dot a_i=-\frac{g_i+\dot\lambda}{A_i},\qquad
 \Psi_u=-\lambda,\quad\Psi_v=\sum_i\phi_v,\quad
 \Psi_{w_i}=\phi_w.
\end{equation}
Once $\lambda$ and the values $a_i$ solving
$\phi_a(a_i,v,w_i)=-\lambda$ are known, these identities evaluate a
Hessian-vector product in $O(r)$ scalar operations. Equivalently, the
Hessian is an arrowhead matrix plus a rank-one correction. This count
excludes the cost of root finding to a given precision; we give no
exact finite-arithmetic root-finding algorithm. Sparse-plus-low-rank
Hessian structure for related nonsymmetric cones is known
\citep{ChenGoulart2025}.

For an integer cap $d\geq3$, set
$k_g=\lceil N/(d-2)\rceil$ and
$k_t=\lceil(N-1)/(d-2)\rceil$. The preceding results give the
following parameters for three lifts of $B_p^N$: one scalar power cone
per coordinate, grouped perspective cones $\mathcal P_{p,r}$, and the
norm-cone tree from the proof of Theorem~\ref{thm:lp-granularity}.
\begin{center}
\begin{tabular}{@{}lcc@{}}
\toprule
Representation of $B_p^N$ & Bounds on ambient $\nu_{\rm LH}$
 & Scalar-power barrier\\
\midrule
Coordinate power cones & $[Nh(p),3N]$ & $3N$\\
Grouped perspectives & $[k_gh(p),N+2k_g]$ & $3N$\\
Norm-cone tree & $[k_th(p),N-1+2k_t]$ & $3(N-1+k_t)$\\
\bottomrule
\end{tabular}
\end{center}
For the last entry, restrict a grouped barrier to $u=v=t$ in every
norm-cone block; the tree has $N-1+k_t$ norm coordinates.
The upper bounds in the middle column use barriers with parameter
equal to the cone dimension, which are only known to exist; the last
column uses the explicit barrier built from
\eqref{eq:chares-scalar-power}. These guarantees concern different
barriers and cannot be combined. At $p=2$ the exact
ambient optima improve to $2N,2k_g,2k_t$.
The generalized power cone theorem of \citet{RoyXiao2022} concerns
Euclidean vector outputs; it gives no barrier of constant parameter for
the $\ell_p$ vector output in
\eqref{eq:grouped-p-cone}.

\subsection{Lower bounds for scaled characteristic barriers}
For a proper cone $C\subset\R^m$, define its characteristic integral
$Z_C(x)=\int_{C^*}\exp(-\ip xy)\,dy$.
Its logarithm is $m$-logarithmically homogeneous and, by the
universal-barrier theorem, a self-concordant barrier with parameter $m$
\citep{LeeYue2021}. Among the multiples $\kappa\log Z_C$, a parameter
below $m$ requires $\kappa<1$, and self-concordance then needs a
separate proof. We call the functions $\kappa\log Z_C$, $\kappa>0$,
scaled characteristic barriers; they form the characteristic family.

\begin{theorem}[Boundary asymptotics of the characteristic family]
\label{thm:characteristic-obstruction}
Let $p>2$ and $r\geq2$. If $\kappa\log Z_C$ is self-concordant, its
degree $\nu$ of logarithmic homogeneity satisfies
\begin{align}
 C=\mathcal P_{p,r}:\quad
 \nu&\geq\frac{r+2}{3/2+(r-1)/p}>2,
 \label{eq:grouped-characteristic-bound}\\
 C=K_{p,r+1}:\quad
 \nu&\geq\frac{r+1}{1+(r-1)/p}>2.
 \label{eq:norm-characteristic-bound}
\end{align}
Both right sides tend to $p$ as $r\to\infty$. For $1<p<2$, the same
bounds, with $q=p/(p-1)>2$ in place of $p$, apply to the Legendre duals
of scaled characteristic barriers of the dual cone.
\end{theorem}
\begin{proof}
Write $\alpha=1/p$, $\beta=1/q$. Parameterize the interior of the dual
in \eqref{eq:perspective-dual} by
\[
 a=\alpha s,\qquad z=-sy,\qquad
 b=s(\beta\norm y_q^q+\rho),\quad s,\rho>0,\ y\in\R^r.
\]
The absolute Jacobian is $\alpha s^{r+1}$. At
$x_\delta=(1+\delta,1,e_1)$, integration first in $\rho$ and then in
$s$ gives
\begin{equation}\label{eq:grouped-characteristic-integral}
 Z_C(x_\delta)=\alpha\Gamma(r+1)
 \int_{\R^r}[\alpha\delta+g(y)]^{-(r+1)}\,dy,
 \qquad g(y)=\alpha+\beta\norm y_q^q-y_1.
\end{equation}
Young's inequality shows that $g$ vanishes only at $e_1$. Near this
point, with $y_1=1+\eta$,
\[
 g(y)=\tfrac{q-1}{2}\eta^2+\beta\norm{y_\perp}_q^q+O(|\eta|^3).
\]
On a sufficiently small fixed neighborhood it is bounded above and
below by positive multiples of $\eta^2+\norm{y_\perp}_q^q$.
On compact sets outside that neighborhood it has a positive minimum,
and for large $y$ it is bounded below by a positive multiple of
$\norm y_q^q$. For derivative orders $j=0,1,2,3$, differentiation
under the integral replaces the exponent $r+1$ by $r+1+j$ and
multiplies by $(-\alpha)^j(r+1)_j$. The outside integrals are bounded
uniformly as $\delta\downarrow0$, since $q(r+1+j)>r$.

Inside the neighborhood substitute
$\eta=\delta^{1/2}t$, $y_\perp=\delta^{1/q}z$.
The resulting integrands are dominated by a constant multiple of
$(1+t^2+\norm z_q^q)^{-(r+1+j)}$, which is integrable: the anisotropic
volume exponent is $a=1/2+(r-1)/q<r+1$.
Dominated convergence gives, with $B=r+1-a=3/2+(r-1)/p$,
\begin{equation}\label{eq:characteristic-differentiated}
 \frac{d^j}{d\delta^j}Z_C(x_\delta)
 \sim(-1)^j(B)_j C\delta^{-B-j},\qquad 0\leq j\leq3,
\end{equation}
for a constant $C>0$. The coefficient ratios follow either by the beta
integral in anisotropic radial coordinates or by differentiating that
homogeneous model integral. Thus $f=\kappa\log Z_C(x_\delta)$ obeys
$f''\sim\kappa B\delta^{-2}$ and
$f'''\sim-2\kappa B\delta^{-3}$. Univariate self-concordance forces
$\kappa B\geq1$. Since $\nu=(r+2)\kappa$, this proves the first bound.

For $K_{p,r+1}$ the dual parameterization $(s,sy)$, $y\in B_q^r$,
has Jacobian $s^r$, so at $\widehat x_\delta=(1+\delta,e_1)$,
\[
 Z_C(\widehat x_\delta)
 =\Gamma(r+1)\int_{B_q^r}(1+\delta+y_1)^{-(r+1)}\,dy.
\]
Writing $a=1+y_1$, the transverse section volume is
$\operatorname{vol}(B_q^{r-1})[1-|1-a|^q]^{(r-1)/q}$,
which equals a positive constant times
$a^{(r-1)/q}(1+O(a))$ near zero. The substitution $a=\delta t$,
with the same differentiated dominated-convergence argument, gives
\eqref{eq:characteristic-differentiated} with
$B=r-(r-1)/q=1+(r-1)/p$. This proves the second bound.
The strict inequalities are equivalent to $(r-1)(1-2/p)>0$.
The final statement follows from \eqref{eq:perspective-dual},
$K_{p,r+1}^*=K_{q,r+1}$, and preservation of parameter by Legendre
duality. The statement concerns Legendre duals because the Legendre
dual of a scaled characteristic barrier need not be one.
\end{proof}

The theorem applies to one family of barriers, not to all LHSC barriers,
so its bounds do not replace the lower bounds in the table.
For $p>2$, the perspective bound is strictly smaller than the norm bound
at the same $r$, since their cross-multiplied difference has the sign
of $(r-1)(1/p-1/2)$. This compares the lower bounds only, not the
optimal parameters.

\subsection{The three-dimensional case and explicit formulas that fail}
For $K_{p,3}$, the general bounds are
\begin{equation}\label{eq:three-dimensional-interval}
 h(p)\leq\nu_{\rm LH}(K_{p,3})\leq3.
\end{equation}
The value is two at $p=2$ and three at the polyhedral endpoints $p=1$
and $p=\infty$ \citep{Hildebrand2013}. We do not determine the optimal
value for other $p$; recent work on barrier design still treats it
as unknown \citep{PirauHildebrand2026}.
For $p>2$, Theorem~\ref{thm:characteristic-obstruction} gives the
stronger bound $\nu\geq3p/(p+1)$ for scaled characteristic barriers.
Chares's numerical experiments suggested that this value suffices; we
do not prove this.

\begin{proposition}[Strict separation from the cross-ratio bound]
\label{prop:characteristic-strict-gap}
For $p>2$, $3p/(p+1)>h(p)$. Consequently no scaled characteristic
barrier on $K_{p,3}$ attains Hildebrand's bound. For $p<2$ the analogous
statement concerns Legendre duals of scaled characteristic barriers on
the dual cone and the bound $3p/(2p-1)$.
\end{proposition}
\begin{proof}
Set $s=p-2>0$ and $A=\gamma^{1/p}$. Equation
\eqref{eq:hildebrand-root} becomes
$A^s(sA+s+1)=1$, whereas
$h(p)=2+sA/(s+1)$ and $3p/(p+1)=2+s/(s+3)$.
The left side is increasing in $A$. At $A_0=(s+1)/(s+3)$ its
logarithm is
$g(s)=\log(2s+3)+(s+1)\log((s+1)/(s+3))$.
Here $g(0)=0$ and, putting $z=2/(s+3)$,
\[
 g'(s)=\log(1-z)+z+\frac{2z}{4-3z}=:H(z),\qquad
 H'(z)=\frac{8-24z+24z^2-9z^3}{(1-z)(4-3z)^2}>0
 \quad(0<z<2/3).
\]
The numerator decreases to zero at $2/3$: its derivative is
$-3(8-16z+9z^2)<0$. Since $H(0)=0$, $g(s)>0$ and $A<A_0$.
The desired inequality follows. Duality gives the last assertion.
\end{proof}

Several explicit barrier forms cannot close the gap in
\eqref{eq:three-dimensional-interval}. First, a hypothetical
$2$-LHSC barrier $F$ on a proper cone must be a Lorentz barrier up to
a linear map and an additive constant. To see this, take $h$ tangent
to a level set of $F$ at $x$, put $a=F''[h,h]$, $b=F'''[h,h,h]$, and
use homogeneity to obtain
\[
 F''[h+tx,h+tx]=a+2t^2,\qquad
 F'''[h+tx,h+tx,h+tx]=b-6at-4t^3.
\]
Self-concordance at $t=\pm\sqrt{a/2}$ forces $b=0$.
For $Q=e^{-F}$ this implies $D^3Q=0$ on triples of vectors tangent to
the level set; homogeneity of degree two gives $D^3Q[x,\cdot,\cdot]=0$.
Hence $Q$ is quadratic. Its Hessian is positive on $x$ and negative
definite on the level-set tangent space, so its signature is
$(1,m-1)$. The barrier boundary
condition identifies the cone with a component of $Q\geq0$.
Thus no non-Lorentz cone has a $2$-LHSC barrier; this alone does not
show that the infimum exceeds two.

For a more concrete obstruction, consider, on $K_{p,3}$,
\[
 F=-\alpha_0\log A-\beta_0\log(A^{2/p}-y^2),\qquad
 A=t^p-|x|^p.
\]
If this is an LHSC barrier, an affine line reaching a generic boundary
point where $A>0$ forces $\beta_0\geq1$ by univariate
self-concordance. Along $(t,x,y)=(1,1-\delta,0)$ the same argument
forces $\alpha_0+2\beta_0/p\geq1$. Its degree of logarithmic homogeneity therefore satisfies
$\nu=p\alpha_0+2\beta_0\geq p>h(p)$ for $p>2$.
For $2<p\leq3$ the nonzero coefficient of $|x|^p$ at $y=0$ also
precludes $C^3$ regularity on the interior axis $x=y=0$.
The still simpler formula
\[
 -a\log(t^p-|x|^p-|y|^p)-b\log t,\qquad a>0,
\]
is not $C^2$ on that axis for $1<p<2$ and has zero transverse Hessian
there for $p>2$.

Adding a regular correction with bounded derivatives cannot repair the
last formula.
On $t=1$, let
$L(x,y)=-a\log(1-|x|^p-|y|^p)$.
For $p\leq3$, $p\ne2$, adding a $C^3$ function cannot cancel its
loss of $C^2$ or $C^3$ regularity on the interior axis.
For $p>3$, put $x_\delta=(1-\delta)^{1/p}$ and
$y_\delta=\delta^{1/(p-2)}$. These points are interior for small
$\delta$, and direct differentiation gives
\[
 L_{yy}(x_\delta,y_\delta)=ap(p-1)+o(1),\qquad
 L_{yyy}(x_\delta,y_\delta)
 =ap(p-1)(p-2)\delta^{-1/(p-2)}(1+o(1)).
\]
If a correction has bounded $yy$ and $yyy$ derivatives along this
sequence, the corrected second derivative stays bounded while the
third diverges. Either convexity fails or self-concordance does.
An admissible correction must therefore be singular in the transverse
direction near the flat boundary; we draw no conclusion about
unrestricted corrections.

Finally, we expand $h(p)$ near the endpoints. Write
$d_0=|1-2/p|$ and $w=W(e^{-1})$, where $W$ is Lambert's function.
Then
\begin{align}
 h(p)&=2+2wd_0+2w^2d_0^2+O(d_0^3),&&p\to2,
 \label{eq:premium-near-two}\\
 h(p)&=3-\frac{\log(2p)+1}{p}
 +\frac{\log^2(2p)-\log(2p)+1}{2p^2}
 +O\left(\frac{\log^3 p}{p^3}\right),&&p\to\infty.
 \label{eq:premium-large-p}
\end{align}
In particular $h(p)=2+w|p-2|+O((p-2)^2)$.
To verify the expansions set $a=\gamma^c$ and $s=1/c-2$.
The root equation is $a^s(1+s(a+1))=1$.
After division of its logarithm by $s$, the analytic implicit-function
theorem at $s=0$ gives
$a=w+\tfrac12w(1+w)s+O(s^2)$, proving
\eqref{eq:premium-near-two}. At $c\downarrow0$, let
$\ell=\log(c/2)$. Substitution in the root equation gives
\[
 \gamma=\frac c2\left[1+\tfrac32c(\ell+1)
                  +O(c^2\ell^2)\right],\qquad
 h(p)=3+c(\ell-1)+\tfrac12c^2(\ell^2+\ell+1)
                  +O(c^3|\ell|^3),
\]
which proves \eqref{eq:premium-large-p}. For $p=1+\varepsilon$,
H\"older duality yields, with $L=\log(2/\varepsilon)$,
$h(p)=3-\varepsilon(L+1)+\tfrac12\varepsilon^2(L^2+L+1)
+O(\varepsilon^3L^3)$.

The scalar cross-ratio bound, the perspective cones, the scalar power
barrier, and the partial-minimization theorem are from the literature.
Our contributions are Lemma~\ref{lem:positive-product} and its
consequences for fixed dictionaries, the lower bounds for scaled
characteristic barriers, and the explicit formulas shown to fail.
To the best of our knowledge, the dimension-dependent bounds
\eqref{eq:grouped-characteristic-bound}--\eqref{eq:norm-characteristic-bound}
have not been stated for these families.
